\section{Exact causal allocation on the circle}\label{sec:circle}
Write $\sigma_1$ for normalized arclength on the unit circle, and set
\begin{equation}\label{eq:circle-beta}
 a_m=\frac{\sin(\pi/m)}{\pi/m},\quad\beta_m=a_m^2\quad(m\ge2),
 \qquad a_1=\beta_1=0.
\end{equation}
These are Euclidean squared-distortion coefficients with centroids allowed in the disk. They are not geodesic distortion coefficients or codepoints constrained to the circle.

\begin{lemma}[Circle centroids]\label{lem:circle}
For every $m\ge1$, $q_m(\sigma_1)=1-\beta_m$. Equal arcs and their centroids attain the value, including randomized Borel encoders in the converse.
\end{lemma}
\begin{proof}
For an encoder cell of mass $p$, allowing a fractional membership weight in $[0,1]$, its first moment has norm at most $\sin(\pi p)/\pi$. To see this, test the moment in its own unit direction. Among functions of prescribed integral $p$, the integral of the cosine is maximized by putting weight one on the centered arc of length $2\pi p$; exchanging any lower cosine value inside with a larger one outside improves the integral. Direct integration gives the bound.

Put $f(p)=\sin^2(\pi p)/(\pi^2p)$ and $f(0)=0$. The total retained energy is at most $\sum_j f(p_j)$. The function is concave on $[0,1/2]$: after writing $z=\pi p$, the sign of its second derivative is that of
$h(z)=z^2\cos(2z)-z\sin(2z)+\sin^2z$;
$h(0)=0$ and $h'(z)=-2z^2\sin(2z)\le0$ there.
If $p>1/2$ is one mass, any other positive mass $q$ is at most $1-p$. For $0<u\le1/2$,
\[
 f'(u)-f'(1-u)
 =\frac{\sin^2(\pi u)}{\pi^2u^2(1-u)^2}
 \bigl[2\pi u(1-u)\cot(\pi u)-(1-2u)\bigr]\ge0.
\]
Indeed $\cos(\pi u)\ge1-2u$ by concavity and $\sin(\pi u)\le\pi u$, so the bracket is nonnegative. Concavity on the half interval gives $f'(q)\ge f'(1-p)\ge f'(p)$. Transferring mass from $p$ to $q$ until the larger mass is at most $1/2$ cannot decrease the sum. Zero masses are handled by continuity. Once all masses lie in $[0,1/2]$, Jensen's inequality gives $\sum_j f(p_j)\le m f(1/m)=\beta_m$ for $m\ge2$. For $m=1$ the only mean is zero. Equal arcs have masses $1/m$ and centroids $a_m e^{2\pi i j/m}$, and attain every bound. Use~\eqref{eq:q-energy} to finish.
\end{proof}
The static circle formula is established in the quantization literature~\cite{RosenblattRoychowdhury}; the proof is included to specify the loss and allow arbitrary fractional cells. Our causal use requires the additional compatibility argument below.

\begin{theorem}[Exact sequential product]\label{thm:circle-product}
For all $n\ge1$ and $m_t\ge1$,
\begin{equation}\label{eq:circle-product}
 \calR_{\rm seq}(n;m_1,\ldots,m_n)
 =B+1-\prod_{t=1}^n\beta_{m_t}.
\end{equation}
A deterministic recursive regular-arc encoder attains it. If all factors are positive, every attaining observer satisfies the local optimality and pooled-centroid conditions of Theorem~\ref{thm:energy} at every positive-mass transition.
\end{theorem}
\begin{proof}
The lower is Theorem~\ref{thm:energy} and Lemma~\ref{lem:circle}. Suppose the old conditional means are $r$ times known unit directions. Given an old label, rotate its direction by the report and assign the result to one of $m_t$ equal arcs. Each outgoing label has probability $1/m_t$ in every incoming row. Its conditional mean is $r a_{m_t}$ times the corresponding new direction, independently of the old label. The row quantization is optimal and the incoming centroids agree exactly, so the pooling defect is zero. Induct from $r=1$ at the known preparation to $r=\prod_t a_{m_t}$. The terminal conditional mean is the legal fixed readout of that label. If a factor is zero, thereafter the mean is zero and the zero-response observer attains the same value. The equality criterion when all factors are positive follows by tracing equality through the nonnegative losses.
\end{proof}

\begin{lemma}[The internal deadline consumes disjoint phases]\label{lem:clock}
Suppose the observer has one reused Borel transition rule, its halt decision is a function of the retained label, and every report has the same Haar law independently of prior reports. If it halts exactly after $n\ge1$ reports almost surely, the sets of labels occurring with positive probability after $0,1,\ldots,n$ reports are pairwise disjoint. If their sizes after the positive times are $m_1,\ldots,m_n$, then
\begin{equation}\label{eq:clock-budget}
 1+\sum_{t=1}^n m_t\le M.
\end{equation}
Conversely any such phase alphabets can be implemented by one autonomous program with exactly that many labels.
\end{lemma}
\begin{proof}
Integrating a report row against Haar gives one time-homogeneous transition matrix on labels; halt labels have no further acquisition. If a nonhalt label occurs at two times $s<t<n$, its future stopping-time distribution, conditional on that label, is the same at both times. It must be concentrated at both $n-s$ and $n-t$, a contradiction. A halt label cannot occur before $n$, and a nonhalt label cannot occur at $n$. Initial randomization only increases the number of possible time-zero labels, which is at least one. This proves~\eqref{eq:clock-budget}. For the converse use the disjoint union of the phase alphabets and a single initial label, with each phase row wired to the next and the last phase marked as halt. The report itself supplies no phase. Known pre-report orthogonal actions preserve the integrated Haar law and do not change the argument.
\end{proof}

\begin{theorem}[Exact total-state law and allocation]\label{thm:clock-circle}
For the circle experiment with an internal exact $n$-report deadline, the legal $M$-label class is nonempty exactly when $M\ge n+1$. In that range write
\[
 K=M-1=nq+r,\qquad q\ge1,\quad 0\le r<n.
\]
Then the infimum over all autonomous Borel randomized programs is attained, and
\begin{equation}\label{eq:clock-circle}
 \calR_{\rm aut}(n,M)=B+1-\beta_q^{\,n-r}\beta_{q+1}^{\,r}.
\end{equation}
In particular, the excess equals one when $n+1\le M\le2n$, and
\begin{equation}\label{eq:clock-rate}
 \calR_{\rm aut}(n,M)-B
 \asymp \min\{1,n^3/(M-1)^2\}.
\end{equation}
More generally, for $\mathbb S^d$,
\begin{equation}\label{eq:clock-sphere}
 \calR_{\rm aut}(n,M)-B
 \asymp_d\min\{1,n^{1+2/d}/(M-1)^{2/d}\}.
\end{equation}
Preparation and audit add two physical calls, so these formulas use $n=N-2$. No free-clock interpretation is made.
\end{theorem}
\begin{proof}
The clock lemma proves infeasibility below $n+1$. It also partitions every legal observer into a sequential architecture satisfying~\eqref{eq:clock-budget}. The product theorem bounds its energy by $\prod_t\beta_{m_t}$ and supplies a program attaining every such allocation. For real $x>1$ the increasing function
$\ell(x)=2\log(\sin(\pi/x)/(\pi/x))$
is strictly concave, because with $z=\pi/x$,
\[
 \ell''(x)=-\frac2{x^2}\bigl[(1-z\cot z)^2+z^2\bigr]<0.
\]
An exchange between two integer widths differing by at least two increases their product. Thus positive products are maximized by the balanced allocation $q,q+1$. If $K<2n$, every allocation has a width one and product zero, so the same formula remains valid without taking $\log\beta_1$. Widths can use all $K$ labels since $\beta_m$ increases. The explicit regular-arc construction proves attainment and the asserted threshold.

The rate~\eqref{eq:clock-rate} follows from $1-\beta_m\asymp m^{-2}$ and the product bounds. For the sphere, Jensen gives $\sum_t m_t^{-2/d}\ge n^{1+2/d}K^{-2/d}$, while the balanced integer allocation has a sum at most $2^{2/d}$ times this quantity. Apply Theorem~\ref{thm:sphere}. These arguments include $q=1$ by the truncation at one.
\end{proof}

\begin{corollary}[Resolution and synthesis]\label{cor:resolution-clock}
For $0<\rho\le\rho_0<1$, excess at most $\rho^2$ on the sphere requires and, up to constants depending only on $d,\rho_0$, is attained by
\[
 M-1\asymp_d n^{1+d/2}\rho^{-d}.
\]
A checkpoint requires $m\asymp_d\rho^{-d}$ instead. On the circle the optimal phase widths require one quotient and remainder computation on the binary integers $M-1,n$; ordinary long division uses $O((\log M+\log n)^2)$ bit operations. This is an exact allocation result, not an efficient global synthesis claim for arbitrary kernels. Compiling all transition rows and their real constants has the separate cost stated in Proposition~\ref{prop:numerics}.
\end{corollary}
\begin{proof}
Solve~\eqref{eq:clock-sphere} in the nonsaturated range, and use the checkpoint bound. The arithmetic claim follows from the balanced formula; enumerating a table of all states still costs at least its output length.
\end{proof}
